\begin{thebibliography}{99}
\bibitem{bib-bib1} 
\par\noindent    Samson Abramsky and Achim Jung.  Handbook of Logic in Computer Science, volume 3, chapter Domain theory, pages 1–168.  Clarendon Press, Oxford, 1994.  

\bibitem{bib-bib2} 
\par\noindent    Jiří Adámek and Jiří Rosickỳ.  Locally presentable and accessible categories, volume 189.  Cambridge University Press, 1994.  

\bibitem{bib-bib3} 
\par\noindent    Gregory Arone and Michael Ching.  Goodwillie calculus.  In Handbook of homotopy theory, pages 1–38. Chapman and Hall/CRC, 2020.  

\bibitem{bib-bib4} 
\par\noindent    Michael Barr.  Coalgebras over a commutative ring.  Journal of Algebra, 32(3):600–610, 1974.  

\bibitem{bib-bib5} 
\par\noindent    Michael Barr.  $*$-autonomous categories and linear logic.  Mathematical Structures in Computer Science, 1(2):159–178, 1991.  doi:10.1017/S0960129500001274.  

\bibitem{bib-bib6} 
\par\noindent    Jon Beck.  Distributive laws.  In Seminar on Triples and Categorical Homology Theory, pages 119–140, Berlin, Heidelberg, 1969. Springer Berlin Heidelberg.  

\bibitem{bib-bib7} 
\par\noindent    Nick Benton.  A mixed linear and non-linear logic: Proofs, terms and models.  In International Workshop on Computer Science Logic, pages 121–135. Springer, 1994.  

\bibitem{bib-bib8} 
\par\noindent    John D. Berman.  Higher lawvere theories.  Journal of Pure and Applied Algebra, 224(9):106362, 2020.  

\bibitem{bib-bib9} 
\par\noindent    Gavin M. Bierman.  What is a categorical model of intuitionistic linear logic?  In Typed Lambda Calculi and Applications, Second International Conference on Typed Lambda Calculi and Applications, TLCA, volume 902 of Lecture Notes in Computer Science, pages 78–93. Springer, 1995.  doi:10.1007/BFB0014046.  

\bibitem{bib-bib10} 
\par\noindent    Pierre-Louis Curien.  Operads, clones, and distributive laws.  In Operads and universal algebra, pages 25–49. World Scientific, 2012.  

\bibitem{bib-bib11} 
\par\noindent    Valeria de Paiva.  Categorical Semantics of Linear Logic for All, pages 181–192.  Springer Netherlands, Dordrecht, 2014.  

\bibitem{bib-bib12} 
\par\noindent    Sanath Devalapurkar and Peter Haine.  On the James and Hilton-Milnor splittings, and the metastable EHP sequence.  Documenta Mathematica, 26:1423–1464, 2021.  

\bibitem{bib-bib13} 
\par\noindent    Eric Finster, Samuel Mimram, Maxime Lucas, and Thomas Seiller.  A cartesian bicategory of polynomial functors in homotopy type theory.  In Proceedings 37th Conference on Mathematical Foundations of Programming Semantics, MFPS 2021, volume 351 of EPTCS, pages 67–83, 2021.  doi:10.4204/EPTCS.351.5.  

\bibitem{bib-bib14} 
\par\noindent    Marcelo Fiore, Zeinab Galal, and Hugo Paquet.  Stabilized profunctors and stable species of structures.  Logical Methods in Computer Science, 20, 2024.  doi:10.46298/LMCS-20(1:17)2024.  

\bibitem{bib-bib15} 
\par\noindent    Marcelo Fiore, Nicola Gambino, and Martin Hyland.  Monoidal bicategories, differential linear logic, and analytic functors.  Preprint, 2024.  arXiv:2405.05774.  

\bibitem{bib-bib16} 
\par\noindent    Marcelo Fiore, Nicola Gambino, Martin Hyland, and Glynn Winskel.  The cartesian closed bicategory of generalised species of structures.  Journal of the London Mathematical Society, 77(1):203–220, 2008.  

\bibitem{bib-bib17} 
\par\noindent    Grant A. Fraser.  The semilattice tensor product of distributive lattices.  Transactions of the American Mathematical Society, 217:183–194, 1976.  

\bibitem{bib-bib18} 
\par\noindent    Nicola Gambino and Joachim Kock.  Polynomial functors and polynomial monads.  Mathematical proceedings of the cambridge philosophical society, 154(1):153–192, 2013.  

\bibitem{bib-bib19} 
\par\noindent    David Gepner, Rune Haugseng, and Joachim Kock.  $\infty$-operads as analytic monads.  International Mathematics Research Notices, 2022(16):12516–12624, 2022.  

\bibitem{bib-bib20} 
\par\noindent    Jean-Yves Girard.  Linear logic.  Theoretical computer science, 50(1):1–101, 1987.  doi:10.1016/0304-3975(87)90045-4.  

\bibitem{bib-bib21} 
\par\noindent    Jean-Yves Girard.  Normal functors, power series and $\lambda$-calculus.  Annals of pure and applied logic, 37(2):129–177, 1988.  doi:10.1016/0168-0072(88)90025-5.  

\bibitem{bib-bib22} 
\par\noindent    Thomas G Goodwillie.  Calculus III: Taylor series.  Geometry \& Topology, 7(2):645–711, 2003.  

\bibitem{bib-bib23} 
\par\noindent    Elies Harington and Samuel Mimram.  Polynomials in homotopy type theory as a Kleisli category.  In Proceedings of MFPS XL, volume 4 of Electronic Notes in Theoretical Informatics and Computer Science, 2024.  

\bibitem{bib-bib24} 
\par\noindent    Rune Haugseng.  On (co) ends in $\infty$-categories.  Journal of Pure and Applied Algebra, 226(2):106819, 2022.  

\bibitem{bib-bib25} 
\par\noindent    Simon Henry and Nicholas J. Meadows.  Higher theories and monads.  Preprint, 2021.  arXiv:2106.02706.  

\bibitem{bib-bib26} 
\par\noindent    Martin Hyland and Andrea Schalk.  Abstract games for linear logic extended abstract.  Electronic Notes in Theoretical Computer Science, 29:127–150, 1999.  

\bibitem{bib-bib27} 
\par\noindent    Martin Hyland and Andrea Schalk.  Glueing and orthogonality for models of linear logic.  Theoretical computer science, 294(1-2):183–231, 2003.  doi:10.1016/S0304-3975(01)00241-9.  

\bibitem{bib-bib28} 
\par\noindent    Bart Jacobs.  Semantics of weakening and contraction.  Annals of Pure and Applied Logic, 69(1):73–106, 1994.  doi:10.1016/0168-0072(94)90020-5.  

\bibitem{bib-bib29} 
\par\noindent    André Joyal.  Quasi-categories and Kan complexes.  Journal of Pure and Applied Algebra, 175(1-3):207–222, 2002.  

\bibitem{bib-bib30} 
\par\noindent    William F. Keigher.  Symmetric monoidal closed categories generated by commutative adjoint monads.  Cahiers de topologie et géométrie différentielle, 19(3):269–293, 1978.  

\bibitem{bib-bib31} 
\par\noindent    Marie Kerjean.  Weak topologies for linear logic.  Logical Methods in Computer Science, 12, 2016.  doi:10.2168/LMCS-12(1:3)2016.  

\bibitem{bib-bib32} 
\par\noindent    Joachim Kock.  Data types with symmetries and polynomial functors over groupoids.  Electronic Notes in Theoretical Computer Science, 286:351–365, 2012.  

\bibitem{bib-bib33} 
\par\noindent    Joachim Kock, André Joyal, Michael Batanin, and Jean-François Mascari.  Polynomial functors and opetopes.  Advances in Mathematics, 224(6):2690–2737, 2010.  

\bibitem{bib-bib34} 
\par\noindent    Yves Lafont.  Logiques, catégories et machines.  PhD thesis, Université Paris 7, 1988.  

\bibitem{bib-bib35} 
\par\noindent    Jacob Lurie.  Higher topos theory.  Princeton University Press, 2009.  

\bibitem{bib-bib36} 
\par\noindent    Jacob Lurie.  Higher algebra.  Preprint, 2017.  URL: https://people.math.harvard.edu/˜lurie/papers/HA.pdf.  

\bibitem{bib-bib37} 
\par\noindent    Jacob Lurie.  Kerodon.  https://kerodon.net, 2018.  

\bibitem{bib-bib38} 
\par\noindent    Paul-André Mellies.  Categorical semantics of linear logic.  Panoramas et syntheses, 27:15–215, 2009.  

\bibitem{bib-bib39} 
\par\noindent    Paul-André Melliès and Nicolas Tabareau.  Resource modalities in tensor logic.  Annals of Pure and Applied Logic, 161(5):632–653, 2010.  doi:10.1016/J.APAL.2009.07.018.  

\bibitem{bib-bib40} 
\par\noindent    Paul-André Melliès, Nicolas Tabareau, and Christine Tasson.  An explicit formula for the free exponential modality of linear logic.  Mathematical Structures in Computer Science, 28(7):1253–1286, 2018.  doi:10.1017/S0960129516000426.  

\bibitem{bib-bib41} 
\par\noindent    Paul-André Melliès.  Categorical models of linear logic revisited.  Unpublished manuscript, 2003.  URL: https://hal.science/hal-00154229.  

\bibitem{bib-bib42} 
\par\noindent    nLab authors.  commutative algebraic theory.  https://ncatlab.org/nlab/show/commutative+algebraic+theory, February 2025.  Revision 39.  

\bibitem{bib-bib43} 
\par\noindent    nLab authors.  Goodwillie calculus.  https://ncatlab.org/nlab/show/Goodwillie+calculus\#GoodwillieDerivativeOfExponentialModality, February 2025.  Revision 68.  

\bibitem{bib-bib44} 
\par\noindent    Emily Riehl.  Categorical homotopy theory, volume 24.  Cambridge University Press, 2014.  

\bibitem{bib-bib45} 
\par\noindent    R. A. G. Seely.  Linear logic, $\ast$-autonomous categories and cofree coalgebras.  In Categories in computer science and logic, Contemp. Math., pages 371–382. American Mathematical Society, 1989.  

\bibitem{bib-bib46} 
\par\noindent    Moss E. Sweedler.  Hopf algebras, volume 44 of Mathematics lecture note series.  W. A. Benjamin, 1969.  

\bibitem{bib-bib47} 
\par\noindent    The Univalent Foundations Program.  Homotopy Type Theory: Univalent Foundations of Mathematics.  https://homotopytypetheory.org/book, Institute for Advanced Study, 2013.  

\bibitem{bib-bib48} 
\par\noindent    Benoît Valiron and Steve Zdancewic.  Modeling simply-typed lambda calculi in the category of finite vector spaces.  Scientific Annals of Computer Science, 24(2), 2014.  doi:10.7561/SACS.2014.2.325.  

\end{thebibliography}
